\chapter{From the page to a distribution over sounds}

This chapter follows one Salishan word from the mark on the page to a symbol a distribution can be
taken over, and says which stage each measurement in this work belongs to. A measurement on this
corpus is only interpretable once it is known what its symbols stand for.

\section{Three stages, and they are not interchangeable}

A symbol in this corpus can mean three different things, and a distribution taken over each of them
answers a different question.

\begin{enumerate}
  \item A \textbf{character} is a codepoint in a file. A distribution over characters is a
        distribution over what a transcriber typed.
  \item A \textbf{segment} is a base with everything that modifies it. A distribution over segments
        is a distribution over the units the orthography writes.
  \item A \textbf{sound} is the referent of the segment, with the definitions of one sound collapsed.
        A distribution over sounds is the first one that is about the language.
\end{enumerate}

Each stage discards something the stage before it carried, and each discard is a claim that has to
be true. Skipping a stage does not leave the measurement neutral. It leaves it measuring the thing
that stage was there to remove.

\section{Stage zero, the page}

Nothing above starts until somebody has read the paper. A text extraction of a scanned typescript is
not the paper: all four ICSNL 2 papers held here \cite{src:Elmendorf-1967, src:Hamp-1967,
src:Hess-1967, src:Kinkade-1967} carry none of their own orthography in their text layer, zero of
the nine marks the corpus uses. They are listed as having no orthography in the text layer, so that
no check grades a correct reading against a file the reading is right to differ from.

The trap runs the other way as well. \emph{The Moon and the Birchbark Canoe} \cite{src:Givens-and-Hall-2023} is born digital and its
text layer is faithful, but the embedded font draws U+026C, the l with a belt, as an l with a bar
through the stem. A transcriber working from page renders sees ł and writes ł. The paper's text
layer holds U+026C 112 times and U+0142 not once. The glyph on the page and the character the page
is drawing are two different questions, and the hand extraction records the second with the first
written down beside it.

\section{Stage one, characters}

The text space is the union of every character these orthographies write their languages with. It
is built from the marked consonants and the schwa, the digit 7 of the practical orthographies, the
stacked glottalization marks, the labialization modifier and Nater's raised space.

Two properties of it bear on every measurement.

\textbf{It is a membership test and not an alphabet.} Its job is to answer whether a token is in one
of these languages. A per-paper set is written as the text space plus what that paper adds, never as
that paper's own alphabet: a test built on one paper's characters reads every other paper as holding
no language, and then reports nothing missing from them either.

\textbf{It carries no phonetic content at all.} \texttt{ɬ} and \texttt{ł} are both in it, as separate
members, because both occur. Nothing in the text space says they are one sound.

Typographic ligatures are repaired before anything else, without a test, because \texttt{fi} is
\texttt{fi} and no decision is involved. Without that repair a case sensitive match counts five and
five as different words.

\section{Stage two, segments}

\texttt{ʷ} is not a segment. It is labialization on the consonant before it. \texttt{̓} is not a
segment either; it is glottalization on the base it sits over. Counting per character measures two
halves of one segment as though they were independent.

The rule is to join every combining mark and every modifier letter to the base in front of it. The
hand extractions already do this when they write a form. Applied to Bev Phillips's
\emph{When Old One Created the Earth} \cite{src:Hall-and-Phillips-2025}, it turns 5552 characters over 35 distinct into 4839 segments
over 50 distinct. The inventory grows under clustering instead of shrinking, because \texttt{kʷ} and
\texttt{k} stop being one letter counted twice. Twenty of the fifty carry a mark or a modifier:

\begin{quote}
\texttt{c̓ kʷ k̓ k̓ʷ m̓ n̓ p̓ qʷ q̓ q̓ʷ w̓ xʷ x̣ x̣ʷ y̓ z̓ ƛ̓ ə́ ʕʷ ʕ̓ʷ}
\end{quote}

\texttt{ə́} is one of the fifty. It is the segment that carries the Lushootseed dialect border. At
the character level it is two codepoints, and the per-run test finds it at width two for that
reason.

\section{Stage three, sounds}

A correspondence table takes a definition and returns the sound it stands for. The definitions that
differ across the corpus, counted by segment over 23 hand extractions:

\tiny
\setlength{\tabcolsep}{2pt}
\begin{longtable}{@{}>{\arraybackslash}p{\dimexpr\textwidth/4-2\tabcolsep\relax}>{\arraybackslash}p{\dimexpr\textwidth*3/4-2\tabcolsep\relax}@{}}
\toprule
Sound & How the corpus writes it \\
\midrule
lateral fricative & \texttt{ɬ} in seventeen papers, \texttt{ł} in four. LaFontaine and Janzen \cite{src:LaFontaine-and-Janzen-2024} write \texttt{ł} 329 times and \texttt{ɬ} never. Nater's voiceless paper writes \texttt{ł} 98 times and \texttt{ɬ} three. \\
glottal stop & \texttt{ʔ} in most, the digit \texttt{7} in the van Eijk orthography. Alexander and Davis \cite{src:Alexander-and-Davis-2026} write \texttt{7} 934 times and \texttt{ʔ} four. Matthewson and Redan \cite{src:Matthewson-and-Redan-2026} write \texttt{7} 125 times and \texttt{ʔ} twice. A capital \texttt{Ɂ} occurs in one paper. \\
uvular fricative & \texttt{χ}, \texttt{x̌} and \texttt{x̣}, all three in six or more papers. \\
schwa & \texttt{ə} at U+0259 and \texttt{ǝ} at U+01DD, a turned e and a different codepoint. \\
\bottomrule
\end{longtable}
\normalsize

Three of those four are named in the references chapter. The fourth comes out of counting, and so
do the capital glottal stop and the aspirated \texttt{pʰ} and \texttt{tʰ} that only Nater writes.
That is the argument for building the table by counting the corpus instead of from a description of
it.

\subsection{Where the correspondences come from}

Nater states his own, in a paper this corpus holds \cite{src:Nater-2024}: \texttt{/ł/ = /ɬ/},
\texttt{/χ/ = /x̌/}, \texttt{/χʷ/ = /x̌ʷ/}, \texttt{/p’, t’, q’/ = /p̓, t̓, q̓/},
\texttt{/kʷ’, qʷ’/ = /k̓ʷ, q̓ʷ/} and \texttt{/’/ = [ʔ]}. A linguist writing down what his own symbols
equal is the strongest source available here. The van Eijk 7 is declared by its orthography. The
rest follow the standard descriptions of Americanist transcription, they are assigned by hand, and
any error in them is an error in every result taken downstream.

\subsection{Two collisions the table must avoid}

\texttt{kʷ} and \texttt{qʷ} must not go to one symbol. A labialized velar and a labialized uvular
are different phonemes in these languages, and the collapse shows itself at once: \texttt{kʷɬ},
\texttt{qɬ} and \texttt{qʷł} arrive as one form, which merges three words.

The uvular fricative must not take the symbol \texttt{x}. A plain velar \texttt{x} passes through as
that same symbol, and the mapped uvular and the unmapped velar then meet in one cell. A pass-through
is not free of the namespace it passes into.

Every sound therefore takes a codepoint of its own from the private use area, assigned by the table
and never typed. No sound can collide with a letter the corpus writes. A table like this is exactly
where a wrong assertion is invisible: nothing downstream can tell a real merger from a typo, and
only reading its output for merged forms finds one.

\subsection{What the mapping does not claim}

Mapping \texttt{ɬ} and \texttt{ł} to one symbol says two definitions are one sound. That is a fact
about transcription practice. Whether Nuxalk's lateral fricative and Lushootseed's are the same
sound is a question for a comparativist, and nothing here asserts it. The mapping buys something
narrow and it is enough: a difference measured afterward is not the difference between two typists.

\section{What a distribution reads at each stage}

\tiny
\setlength{\tabcolsep}{2pt}
\begin{longtable}{@{}>{\arraybackslash}p{\dimexpr\textwidth/4-2\tabcolsep\relax}>{\arraybackslash}p{\dimexpr\textwidth*3/4-2\tabcolsep\relax}@{}}
\toprule
Over what & What the number is about \\
\midrule
bytes & the encoding, and whether a mark is precomposed or decomposed \\
characters & the transcriber, and which conference volume the paper appeared in \\
segments & the orthography, comparable inside one paper and across papers sharing a convention \\
sounds & the language, comparable across every paper in the corpus \\
\bottomrule
\end{longtable}
\normalsize

The failure this ordering prevents has a definite shape. A model counting character bigrams reads
two dialects that share a word as maximally far apart, because it is measuring the transcriber and
not the speaker. The distance is real, it is stable, it will replicate, and it is about the wrong
thing.

\section{Limitations}

\textbf{The correspondence table is not part of the screening.} The measurements of the other
chapters take their distributions over bytes, and none of them is taken over sounds.

\textbf{The clustering rule has a single implementation.} Stage two is applied in one experiment and
is not shared by the other measurements.

\textbf{The word web is built on written forms.} Nothing takes a character out of the text space and
hands back what it stands for at the point the web is built. Two dialects writing one word two ways
are still two nodes. Stage three would close it.
